\section{Normal coefficients at multiple zeros and simple tangencies}
\label{sec:motion-bending}

The denominator translation has an explicit second jet. We retain
that jet because the first velocity can be vertical. Throughout
this section, a normal coefficient means the real part of a
coefficient in the coordinate $s-\rho$, with the positive real
parameter oriented by
\[
 F(s,u)=\sum_{n\ge1}\frac{\chi(n)}{(n+u)^s}
 =L(s)-usL(s+1)+\frac{u^2s(s+1)}2L(s+2)+O(u^3).
\]
The conductor-one case is $L=\zeta$. The equality first holds
in the absolutely convergent half-plane; its normally convergent
rising-factorial expansion continues it to a neighborhood of
each critical-line zero. The Euler product, primitive functional
equation and endpoint nonvanishing used below have the exact
normalizations in \cite[Section~25.15, equations (2), (5), (6), (9)]{dlmf}.

\begin{lemma}[Reflected logarithmic derivative]
\label{lem:reflected-positive-new}
Let $\Lambda$ be an entire function of order at most one, with
infinitely many zeros in $0\le\Re s\le1$, satisfying
$\Lambda(s)=\epsilon\overline{\Lambda(1-\bar s)}$, $|\epsilon|=1$.
At a zero $\rho=1/2+i\gamma$ of multiplicity $m\ge1$,
\begin{equation}\label{eq:reflected-positive-new}
 \Re\frac{\Lambda'}{\Lambda}(\rho+1)
 =m+\sum_{\rho'\ne\rho}
 \frac{3/2-\Re\rho'}{(3/2-\Re\rho')^2+(\gamma-\Im\rho')^2}>m.
\end{equation}
The sum counts multiplicity and converges absolutely.
\end{lemma}

\begin{proof}
Set $G(z)=\Lambda(\rho+z)$ and write
$G(z)=b_mz^m+b_{m+1}z^{m+1}+\cdots$. Reflection gives
$\Re(b_{m+1}/b_m)=0$. Hadamard factorization yields
\[
 G(z)=z^m e^{az+b}\prod_{\alpha\ne0}
       (1-z/\alpha)e^{z/\alpha},\qquad
 \alpha=\rho'-\rho,\quad a=b_{m+1}/b_m.
\]
Here $\sum|\alpha|^{-2}<\infty$. At $z=1$, the real part
of the logarithmic derivative is
$m+\sum\Re((1-\alpha)^{-1}+\alpha^{-1})$.
Since $|\Re\alpha|\le1/2$, the real parts of
$\alpha^{-1}$ are absolutely summable. Reflection pairs
$\alpha$ with $-\bar\alpha$, so their sum is zero.
Each remaining term is positive, and its tail is
$O(|\alpha|^{-2})$. This proves the formula and its strict sign.
\end{proof}

For a nonprincipal primitive character of parity $a\in\{0,1\}$,
put
\[
 Q_\chi(s)=(q/\pi)^{(s+a)/2}\Gamma((s+a)/2),\qquad
 \Lambda_\chi=Q_\chi L.
\]
The theta integral proves that $\Lambda_\chi$ is entire of
order at most one: after splitting its Mellin integral at one,
Poisson summation replaces the small-variable part by a conjugate
theta series that decays exponentially. The resulting integral
has growth $\exp(O_q(R\log(R+2)))$ on $|s|\le R$.
The functional equation cancels each trivial zero against its
gamma pole; at the endpoint this uses $L(1,\chi)\ne0$.
Thus its remaining zeros are precisely the nontrivial ones,
and the reflected lemma applies also for complex $\chi$.

\begin{proposition}[Exact second normal coefficient]
\label{prop:motion-second-jet}
At a critical-line zero of $L$ of order $m\ge1$, write
\[
 L(\rho+w)=A_mw^m+A_{m+1}w^{m+1}+O(w^{m+2}),\qquad
 D_\rho=\frac1\rho+\frac{L'}L(\rho+1)-\frac{A_{m+1}}{A_m}.
\]
For every nonprincipal primitive character,
\begin{equation}\label{eq:dirichlet-bending-positive-new}
 \Re D_\rho>m+h_a(\gamma)>m,
 \qquad
 \begin{cases}
 h_0(\gamma)=\dfrac1{2(\gamma^2+1/4)}
                  -\dfrac\pi{2\cosh(\pi\gamma)},\\[4pt]
 h_1(\gamma)=\dfrac\pi{2\cosh(\pi\gamma)}.
 \end{cases}
\end{equation}
For $L=\zeta$, the corresponding bound is
\begin{equation}\label{eq:zeta-bending-positive-new}
 \Re D_\rho>m-\frac3{2(\gamma^2+9/4)}
                    -\frac\pi{2\cosh(\pi\gamma)}>0.
\end{equation}
Every multiple zero ($m\ge2$) has convergent local sections
after $u=\tau^m$ of the form
\begin{equation}\label{eq:multiple-bending-jets-new}
 s_j(u)=\rho+c_j u^{1/m}
              +\frac{D_\rho}{m}c_j^2u^{2/m}+O(u^{3/m}),
 \qquad c_j^m=\frac{\rho L(\rho+1)}{A_m}\ne0.
\end{equation}
A vertical first coefficient $c_j=i\beta_j$ has strictly
negative second normal coefficient
$-\beta_j^2\Re D_\rho/m$. Consequently every branch has
a nonzero normal coefficient by the second Puiseux order.
When vertical first coefficients occur, the numbers on the
left and right are respectively $(m/2+1,m/2-1)$ for even $m$,
and $((m+1)/2,(m-1)/2)$ for odd $m$.
\end{proposition}

\begin{proof}
Reflection of the completed germ gives
\[
 \Re\frac{A_{m+1}}{A_m}=-\Re\frac{Q_\chi'}{Q_\chi}(\rho).
\]
Substitution expresses $\Re D_\rho$ as the left side of
\eqref{eq:reflected-positive-new} plus
$\Re(1/\rho+Q_\chi'/Q_\chi(\rho)-Q_\chi'/Q_\chi(\rho+1))$.
The conductor cancels. Digamma reflection and recurrence
\cite[Section~5.5, equations (2), (4)]{dlmf} evaluate this last
term as $h_a(\gamma)$. For even parity use
\[
 \Re\{\psi(\rho/2)-\psi((\rho+1)/2)\}
       =-\frac\pi{\cosh(\pi\gamma)};
\]
the recurrence $\psi(1+\rho/2)=\psi(\rho/2)+2/\rho$
gives the odd case. Both corrections are positive, since
$\cosh(\pi\gamma)\ge1+\pi^2\gamma^2/2>
\pi(\gamma^2+1/4)$.

For zeta use the entire completion
$\xi(s)=s(s-1)\pi^{-s/2}\Gamma(s/2)\zeta(s)/2$.
Its pole-removal factors contribute
\[
 \Re\{1/\rho+1/(\rho-1)-1/(\rho+1)\}
             =-3/[2(\gamma^2+9/4)],
\]
giving the first inequality in
\eqref{eq:zeta-bending-positive-new}. For the second inequality we use the following
small-height certificate.
The positive theta representation and the normalization
$H_0(x)=\xi(1/2+ix/2)/8$ give
\[
 \frac{H_0(x)}{H_0(0)}\ge1-c_1x^2,\qquad
 c_1=\frac{\xi''(1/2)}{8\xi(1/2)}.
\]
Ball power-series arithmetic at 80 and 120 decimal digits,
using the regular expression
$\xi(s)=(s-1)\pi^{-s/2}\Gamma(1+s/2)\zeta(s)$,
certifies
\[
 0.00577624827885474269<c_1<0.00577624827885474270<1/128.
\]
Thus $H_0(x)>H_0(0)/2$ for $|x|\le8$, and every
critical-line zeta zero has $|\gamma|>4$. The two subtracted
terms in \eqref{eq:zeta-bending-positive-new} are then less
than $(3/2+1/\pi)/16<11/96<1$. The certificate consists
only of finite ball inequalities; the passage to all these
heights uses $\cos v\ge1-v^2/2$ in the positive integral.

Finally, substitute $w=c\tau+d\tau^2+\cdots$, $u=\tau^m$,
in the displayed deformation. The coefficients of $\tau^m$
and $\tau^{m+1}$ give
\[
 A_mc^m=\rho L(\rho+1),\qquad
 mA_mc^{m-1}d+A_{m+1}c^{m+1}
       -\{L(\rho+1)+\rho L'(\rho+1)\}c=0.
\]
The $u^2$ term has degree $2m\ge m+2$ and does not enter.
The distinct nonzero roots $c$ and the implicit function
theorem give all the sections and their convergent expansions.
Their regular polygon has two vertical vertices for even $m$
and one for odd $m$ whenever it has any; the other vertices
split equally by the sign of their real part. Adding the
strictly left-bending vertices gives the asserted counts.
\end{proof}

\begin{lemma}[Exterior phase bound]
\label{lem:exterior-sector-new}
For every Dirichlet character, with the continuous Euler
argument,
\[
 \left|\arg\frac{L(5/2+it)}{L(3/2+it)}\right|<3/4.
\]
Writing $R=L(\rho+2)/L(\rho+1)$, one consequently has
\begin{equation}\label{eq:exterior-positive-new}
 P_\rho:=\Im((\rho+1)R)
       >|R|\frac{23\gamma-36}{32}>0\qquad(\gamma\ge2).
\end{equation}
\end{lemma}

\begin{proof}
The absolute Euler logarithm is at most
$\log(\zeta(3/2)/\zeta(5/2))$. The integral tail after five,
and the first four terms respectively, give the rational bounds
\[
 \begin{split}
 \zeta(3/2)&<\frac98+\frac5{14}+\frac7{36}+\frac{99}{100}
                =\frac83-\frac1{12600},\\
 \zeta(5/2)&>\frac{33}{32}+\frac7{40}+\frac5{78}
                >\frac{127}{100}.
 \end{split}
\]
These follow from $\sqrt2>7/5$, $\sqrt3>12/7$,
$\sqrt5>20/9$, $\sqrt2<10/7$, and $\sqrt3<26/15$,
each checked by squaring. Also
$e^{3/4}>\sum_{j=0}^4(3/4)^j/j!=4331/2048>800/381$.
This proves the phase bound. For $|\theta|<3/4$,
$\cos\theta>23/32$ and $\sin\theta>-3/4$;
expanding $\Im((3/2+i\gamma)|R|e^{i\theta})$
proves \eqref{eq:exterior-positive-new}.
\end{proof}

\begin{proposition}[The second jet at a simple tangent]
\label{prop:simple-tangent-new}
At every simple critical-line zero, let
$v=\rho L(\rho+1)/L'(\rho)$. The simple section has
\begin{equation}\label{eq:simple-bending-jet-new}
 s(u)=\rho+vu+
 \left(D_\rho v^2-\frac{(\rho+1)L(\rho+2)}{2L(\rho+1)}v\right)u^2
 +O(u^3).
\end{equation}
At a vertical velocity $v=i\beta$, put $d=\Re D_\rho$.
For a primitive nonprincipal character with $\gamma\ge2$,
and for every positive-height zeta zero, $d>0$, $P_\rho>0$,
and the second normal coefficient is exactly
\[
       -\beta^2d+\frac{\beta P_\rho}{2}
       =\beta d\left(\frac{P_\rho}{2d}-\beta\right).
\]
It is strictly negative for $\beta<0$. For $\beta>0$ it is
positive below $P_\rho/(2d)$ and negative above that value;
it vanishes at equality.
\end{proposition}

\begin{proof}
For $m=1$ the $u^2$ deformation term must be retained.
Substitute $s=\rho+vu+bu^2$ and match the coefficient of
$u^2$:
\[
 L'(\rho)b+\tfrac12L''(\rho)v^2
 -\{L(\rho+1)+\rho L'(\rho+1)\}v
 +\tfrac12\rho(\rho+1)L(\rho+2)=0.
\]
Division by $L'(\rho)$ proves the formula. The previous
proposition proves $d>0$ also at multiplicity one, and the
exterior phase lemma proves $P_\rho>0$ in the stated ranges.
Taking real parts with $v=i\beta$ proves every sign assertion.
In the notation of the corrected phase, differentiation gives
$d=\psi_\chi'(\gamma)$ exactly: the completed local germ
contributes $\theta_\chi'$, the exterior logarithm contributes
$\Re(L'/L)(\rho+1)$, and
$(-\arctan(1/(2\gamma)))'=\Re(1/\rho)$.
\end{proof}

\paragraph{The central real-character specialization.}
For a real nonprincipal primitive character, the completed
functional equation has root number $+1$; its exact sign is
proved through the quotient of the quadratic-field and rational
zeta completions in \cite[Section~1, equations (1.3)--(1.6)]{gross-signs}.
The central multiplicity $m$ is therefore even. The cited
zero-location theorem and positivity of the Euler product on
$(1,\infty)$ make the leading central coefficient $A_m$
positive, including $m=0$. The paired product gives
\[
 C_\chi:=2+\frac{L'}L(3/2)+\frac{Q_\chi'}{Q_\chi}(1/2)
 =m+\sum_{\gamma>0}\frac2{1+\gamma^2}
 +\begin{cases}2-\pi/2,&a=0,\\ \pi/2,&a=1.\end{cases}
\]
The sum includes multiplicity. At a central zero,
\eqref{eq:multiple-bending-jets-new} gives
\[ c_j^m=L(3/2)/(2A_m)>0,\qquad D_\rho=C_\chi. \]
For $m\equiv2\pmod4$ its branches split equally left and
right; for $m\equiv0\pmod4$, $m\ge4$, the two vertical
vertices give $(m/2+1,m/2-1)$. These formulas assign the
local branches of the central divisor without asserting that
this divisor has positive multiplicity for any character.

\paragraph{The exact Dedekind map.}
For the norm shift $\sum_{\mathfrak a}(N\mathfrak a+u)^{-s}$,
the same rising-factorial expansion replaces $L$ by $\zeta_K$.
The entire completion has pole-removal factors $s(s-1)$,
$r_1$ gamma factors $\Gamma(s/2)$ and $r_2$ factors $\Gamma(s)$.
The preceding calculation therefore gives the explicit identity
\[
 \Re D_\rho=\Re\frac{\xi_K'}{\xi_K}(\rho+1)
 -\frac3{2(\gamma^2+9/4)}
 -\frac{r_1\pi}{2\cosh(\pi\gamma)}
 -\frac{r_2}{2(\gamma^2+1/4)}.
\]
Its first term exceeds $m$, so $\Re D_\rho>0$ whenever
$\gamma^2\ge ((r_2+3)/2+r_1/\pi)/m$.
For a quadratic field the exact factorization
$\zeta_K=\zeta L(\,\cdot\,,\chi_K)$ gives an exterior argument
of absolute value below $3/2$. The alternating cosine
polynomial gives $\cos(3/2)>7/100$, whence
$P_\rho>|R|(7\gamma/100-3/2)>0$ for $\gamma\ge24$.
Thus the same second-jet sign formula holds there. The displayed completion and the quadratic factorization
give the corresponding maps for these coefficient identities.
